\section{几何光学}

\subsection{光纤}

\begin{proset}
    由透明材料制成的正方体\(abcd\)的截面如图所示。一束单色光从\(ab\)边上的\(O\)点以入射角\(\theta =\) \,\qty{45}{\degree}射入正方体，恰好在\(ad\)边上发生全反射，最后从\(cd\)边的中点射出正方体。已知正方体的边长为\(2l\)，光在真空中的传播速度为\(c\)，不考虑多次反射。求：
    \begin{enumerate}[(1)]
        \item 该透明材料的折射率\(n\)；
        \item 单色光在正方体中传播时间。
    \end{enumerate}
    \begin{tikzpicture}[decoration={markings,mark=at position 0.5 with {\arrow{>}}},font=\scriptsize]
        \draw (0,0) rectangle (2,2);
        \coordinate (o) at(0,{3-sqrt(2)});
        \coordinate (p) at({2-sqrt(2)},2);
        \draw[postaction={decorate},purple] ($(o)+(45:-1)$)--(o)node[below right]{\(O\)};
        \draw[dashed,cyan] ($(o)+(-.8,0)$)--++(1.6,0);
        \node[above] at(0,2){\(a\)};
        \node[left] at(0,0){\(b\)};
        \node[right] at(2,0){\(c\)};
        \node[right] at(2,2){\(d\)};
    \end{tikzpicture}
\end{proset}

\begin{answer}
    （1）\(n = \frac{\sqrt{6}}{2} \)（2）\(t = \frac{3l}{v}\)

    \begin{tikzpicture}[decoration={markings,mark=at position 0.5 with {\arrow{>}}},font=\scriptsize]
        \draw (0,0) rectangle (2,2);
        \coordinate (o) at(0,{3-sqrt(2)});
        \coordinate (p) at({2-sqrt(2)},2);
        \draw[postaction={decorate},purple] ($(o)+(45:-1)$)--(o)node[below right,black]{\(O\)};
        \draw[postaction={decorate},purple] (o)--(p);
        \draw[postaction={decorate},purple] (p)--(2,1);
        \draw[dashed,cyan] ($(o)+(-.8,0)$)--++(1.6,0) ($(p)+(0,-.5)$)--++(0,1);
        \draw[dashed] (p)--(2,3)--(2,2);
        \node[above] at(0,2){\(a\)};
        \node[left] at(0,0){\(b\)};
        \node[right] at(2,0){\(c\)};
        \node[right] at(2,2){\(d\)};
    \end{tikzpicture}
\end{answer}